% !TeX root = ../../Book1.tex
% 纲要标题：### 12.3 不可约性判据
% 纲要位置：计划纲要.md，第 456 行。

\section{不可约性判据}
\label{b1:sec:1203}

% 纲要内容（仅作写作计划，尚非教材正文）：
%
% * Eisenstein 判据及平移使用
% * 模素数约化法
% * 有理根检验的适用范围与失效例子
%

唯一分解定理保证因子分解存在而且唯一，却没有立即告诉我们某个给定多项式是否已经不可约。本节给出几种可以实际使用的判据。它们都只能在相应假设下推出结论；某一种检验没有奏效，并不等于多项式可约。

\subsection{Eisenstein 判据及平移}
\label{b1:subsec:1203:01}

设 \(R\) 为唯一分解整环，\(p\in R\) 为素元。若一个多项式模 \(p\) 后只剩最高次项，那么它的非平凡因子模 \(p\) 后也受到很强的限制。再观察常数项中 \(p\) 的幂次，就得到下面的判据。

\begin{theorem}[Eisenstein]\label{b1:thm:eisenstein}
设 \(f=a_nx^n+\cdots+a_1x+a_0\in R[x]\) 本原，\(n\ge1\)。若存在素元 \(p\)，满足
\[
p\nmid a_n,\qquad p\mid a_i\ (0\le i<n),\qquad p^2\nmid a_0,
\]
则 \(f\) 在 \(R[x]\) 及 \(\operatorname{Frac}(R)[x]\) 中均不可约。
\end{theorem}
\begin{proof}
由\cref{b1:thm:gauss-descent}，只需排除 \(R[x]\) 中两个正次数因子的分解。假设 \(f=gh\)，\(\deg g=r>0\)、\(\deg h=s>0\)，则 \(r+s=n\)。首项系数的积为 \(a_n\)，不被 \(p\) 整除，所以 \(g,h\) 的首项系数模 \(p\) 都非零；约化后次数仍为 \(r,s\)。

在整环 \(R/(p)\) 上，有 \(\bar g\bar h=\bar a_nx^n\)。设 \(\bar g,\bar h\) 中最低非零项的指数分别为 \(u,v\)。最低项系数的积非零，故乘积的最低指数为 \(u+v=n\)。另一方面 \(u\le r,v\le s,r+s=n\)，因此只能有 \(u=r,v=s\)。两者均为正数，所以 \(g,h\) 的常数项都被 \(p\) 整除。这使 \(a_0=g(0)h(0)\) 被 \(p^2\) 整除，与假设矛盾。
\end{proof}

这称为\term[eisenstein-panju]{Eisenstein 判据}[Eisenstein criterion]。例如 \(x^5-6x+3\) 对素数三满足该判据，故在 \(\mathbb Q[x]\) 中不可约。没有出现的系数就是零，当然也被三整除；常数项三不被九整除，是不能遗漏的条件。

\ProseParagraphBreak
有时原多项式的系数没有这种形式，平移后却有。对 \(a\in R\)，代入映射 \(T_a:R[x]\rightarrow R[x]\)、\(f(x)\mapsto f(x+a)\) 保持和、积与单位，而且 \(T_{-a}T_a=T_aT_{-a}=1\)。它是环自同构，因而保持单位与不可约性。这不是仅仅改变某个函数的图像，而是在多项式环内部作可逆代换。

\ProseParagraphBreak
以素数 \(p\) 为例，考虑
\[
\Phi_p(x)=1+x+\cdots+x^{p-1}.
\]
记号 \(\Phi_p\) 的一般推广将在后文介绍；这里仅指所列多项式。由几何级数恒等式，
\[
\Phi_p(x+1)=\frac{(x+1)^p-1}{x}
=\sum_{j=1}^p\binom pj x^{j-1}.
\]
右端说明商确为整数系数多项式。它的首项系数为一，常数项为 \(p\)；其余二项式系数都被 \(p\) 整除，因为 \(1\le j<p\) 时 \(j!\) 与 \((p-j)!\) 均不含素因子 \(p\)，而 \(p!\) 含有这个因子。故 \(\Phi_p(x+1)\) 对 \(p\) 满足 Eisenstein 判据，\(\Phi_p(x)\) 也在 \(\mathbb Q[x]\) 中不可约。

\subsection{模素数约化法}
\label{b1:subsec:1203:02}

把系数约化到有限域，可以把有理数域中的问题转成有限次的尝试。这个办法需要保住原多项式的次数。

\begin{proposition}\label{b1:prop:reduction-irreducible}
设 \(f\in\mathbb Z[x]\) 本原，次数为正，\(p\) 是不整除其首项系数的素数。若逐系数约化所得的 \(\bar f\in\mathbb F_p[x]\) 不可约，则 \(f\) 在 \(\mathbb Q[x]\) 及 \(\mathbb Z[x]\) 中都不可约。
\end{proposition}
\begin{proof}
若 \(f\) 在 \(\mathbb Q[x]\) 中可约，\cref{b1:thm:gauss-descent}给出 \(f=gh\)，其中 \(g,h\in\mathbb Z[x]\) 的次数均为正。由于首项系数的积不被 \(p\) 整除，每个因子的首项系数也不被 \(p\) 整除，所以 \(\deg\bar g=\deg g>0\)、\(\deg\bar h=\deg h>0\)。约化后 \(\bar f=\bar g\bar h\) 就是非平凡分解，与假设矛盾。再由同一个 Gauss 分解定理，得到 \(\mathbb Z[x]\) 中的结论。
\end{proof}

上述方法称为\term[mo-sushu-yuehua-fa]{模素数约化法}[reduction modulo a prime]。例如 \(f=x^3+x+1\) 模二后，在 \(\bar0,\bar1\) 处的值都为 \(\bar1\)，所以没有根。域上的二次或三次多项式若可约，两个正次数因子的次数之和至多为三，其中必有一次因子；由\cref{b1:prop:factor-root}，这等价于存在域内的根。因此 \(\bar f\) 不可约，所给整数多项式在 \(\mathbb Q[x]\) 中不可约。

\ProseParagraphBreak
若约化后可约，只能说这次检验没有得出结论。例如 \(x^2+1\) 在 \(\mathbb Q[x]\) 中不可约，因为有理数的平方不可能等于负一；但它模二后等于 \((x+1)^2\)。首项系数条件也不能删去：
\[
2x^2+3x+1=(2x+1)(x+1)
\]
模二后却成为一次不可约多项式 \(x+1\)。失去的因子被约化成了常数，原来的次数分配便无法恢复。

\subsection{有理根检验的适用范围}
\label{b1:subsec:1203:03}

\begin{proposition}[有理根检验]\label{b1:prop:rational-root}
设 \(f=a_nx^n+\cdots+a_0\in\mathbb Z[x]\)，\(a_n\ne0\)。若非零有理数 \(r/s\) 是 \(f\) 的根，其中 \(r,s\in\mathbb Z\)、\(s>0\)、\(\gcd(r,s)=1\)，则 \(r\mid a_0\)、\(s\mid a_n\)。特别地，首一整系数多项式的每个有理根都是整数。
\end{proposition}
\begin{proof}
把 \(f(r/s)=0\) 乘以 \(s^n\)，得到
\[
a_nr^n+a_{n-1}r^{n-1}s+\cdots+a_1rs^{n-1}+a_0s^n=0.
\]
除最后一项以外，每项都被 \(r\) 整除，因此 \(r\mid a_0s^n\)。由于 \(r\) 与 \(s^n\) 互素，由整数的素因子分解可消去 \(s^n\)，得到 \(r\mid a_0\)。同样，除第一项以外，每项都被 \(s\) 整除，故 \(s\mid a_nr^n\)；互素性给出 \(s\mid a_n\)。若 \(a_n=1\)，正分母只能为一。
\end{proof}

这称为\term[youli-gen-jianyan]{有理根检验}[rational root test]。零是否为根，直接查看 \(a_0\) 是否为零即可。当 \(a_0\ne0\) 时，定理只留下有限个可能的分子和分母，可以逐一代入。例如 \(x^3-2\) 的候选有理根只有 \(\pm1,\pm2\)，代入均不为零；次数为三，故它在 \(\mathbb Q[x]\) 中不可约。这里当然也可以直接使用素数二的 Eisenstein 判据。

\ProseParagraphBreak
次数达到四以后，没有有理根只能排除一次因子，不能排除两个二次因子的乘积。例如
\[
x^4+x^2+1=(x^2+x+1)(x^2-x+1).
\]
左边对任何实数都严格为正，因而没有有理根，却在 \(\mathbb Q[x]\) 中可约。选择判据时，先问它排除了哪些可能的因子次数，往往比机械地尝试更多代入值更有效。
